\documentclass[11pt,a4paper]{article}

\usepackage[margin=2.25cm,headheight=15pt]{geometry}
\usepackage[T1]{fontenc}
\usepackage[utf8]{inputenc}
\usepackage{lmodern}
\usepackage{amsmath,amssymb,amsthm}
\usepackage{booktabs,tabularx,longtable,array}
\usepackage{microtype}
\usepackage{enumitem}
\usepackage{listings}
\usepackage{xcolor}
\usepackage{hyperref}
\usepackage{fancyhdr}
\usepackage{caption}
\usepackage{float}
\usepackage{url}
\usepackage{needspace}

\definecolor{linkblue}{RGB}{17,67,111}
\definecolor{softgray}{RGB}{85,95,106}
\hypersetup{colorlinks=true,linkcolor=linkblue,urlcolor=linkblue,citecolor=linkblue,pdftitle={Reproducible Four-Round Midori-64 Differential Characteristic Analysis},pdfauthor={Egemen Ceken}}
\setlength{\parindent}{0pt}
\setlength{\parskip}{5.2pt}
\setlist[itemize]{topsep=3pt,itemsep=2pt,leftmargin=1.5em}
\pagestyle{fancy}
\fancyhf{}
\fancyhead[L]{\small Reproducible Four-Round Midori-64 Analysis}
\fancyhead[R]{\small Technical Report}
\fancyfoot[C]{\small \thepage}
\renewcommand{\headrulewidth}{0.35pt}
\captionsetup{font=small,labelfont=bf}
\lstset{basicstyle=\ttfamily\footnotesize,breaklines=true,frame=single,
  columns=fullflexible,keepspaces=true,showstringspaces=false,
  rulecolor=\color{black!18},backgroundcolor=\color{black!2},
  xleftmargin=4pt,xrightmargin=4pt,aboveskip=8pt,belowskip=8pt}
\newtheorem{proposition}{Proposition}

\title{\vspace{-1.2cm}\bfseries Reproducible Four-Round Midori-64\\[4pt]\large Differential Characteristic Analysis}
\author{Egemen \c{C}eken\\\small Technical Research Report}
\date{October 2026}

\begin{document}
\maketitle
\vspace{-1.1cm}
\begin{abstract}
This technical report presents a reproducible mixed-integer linear programming (MILP) analysis of four-round
Midori-64. The S-box activity and exact differential transitions, the ShuffleCell
permutation, and the MixColumn layer are encoded using a single consistent modeling
framework adapted from \c{I}lter and Sel\c{c}uk. The complete four-round model has
1,280 binary variables and 2,049 constraints. An exhaustive component check verifies
that the reduced 19-inequality S-box formulation accepts exactly the 97 valid
input/output/probability patterns among all 1,024 binary patterns.
Gurobi 13.0.3 \emph{certifies} that the minimum number of active S-boxes is 16,
with an objective value and best bound both equal to 16 and a zero optimality gap.
An independent word-level dynamic program reproduces the same four-round bound without MILP. A feasible differential characteristic with 16 active S-boxes attains the
DDT-product probability $2^{-32}$; combining this witness with the certified
active-S-box bound proves its optimality \emph{within the specified four-round
characteristic model}. Finally, the characteristic is realized by an explicit
master key and plaintext pair in a four-round reduced Midori-64 implementation
with key addition, round constants, and whitening. We distinguish a
characteristic probability product from an exact fixed-key differential probability
and document exploratory conditional fixed-key sampling near $2^{-22}$ for the exhibited master key. This reproduces known fixed-key phenomena rather than claiming a new security bound for the full cipher.
\end{abstract}

\section{Introduction and Scope}
Midori is a family of lightweight block ciphers designed for low energy
consumption \cite{banik2015}. This study focuses on the 64-bit variant,
Midori-64, whose state consists of sixteen four-bit cells. Our goal is to reproduce the four-round active-S-box bound using an exact MILP model,
derive an optimal DDT-product characteristic within that model, and independently
validate its four-round realization under a concrete master key.

The modeling methodology follows the S-box, permutation, and XOR construction
in \c{I}lter and Sel\c{c}uk's analysis of the \emph{FUTURE} block cipher
\cite{ilter2023}; we apply the \emph{method}, not FUTURE's specific S-box or linear
layer, to Midori-64. Foundational MILP studies include Mouha et al.~\cite{mouha2012} and the convex-hull formulation of Sun et al.~\cite{sun2014}. All component constraints are retained unchanged between the active-S-box
and probability optimizations; only the objective function is changed.

We model four consecutive SubCell layers, indexed $r=0,1,2,3$, with three
intervening ShuffleCell/MixColumn layers. The last round has no following
ShuffleCell/MixColumn operation. This convention also defines the four-round
reduced cipher used in the verification experiment. All cell indices below
follow a column-major order: cells $0,1,2,3$ form the first column, and so on.

\section{Preliminaries}
\subsection{The Midori-64 round components}
The four-bit Midori-64 substitution used in this implementation is
\[
S=(\mathtt{C},\mathtt{A},\mathtt{D},\mathtt{3},\mathtt{E},\mathtt{B},
\mathtt{F},\mathtt{7},\mathtt{8},\mathtt{9},\mathtt{1},\mathtt{5},
\mathtt{0},\mathtt{2},\mathtt{4},\mathtt{6}),
\]
where entries are listed in input order $0,1,\ldots,15$ (hexadecimal notation).
Each state consists of 16 nibbles and thus 64 bits. SubCell applies $S$ to
every nibble. ShuffleCell permutes cell positions, and MixColumn combines the
four cells of each column by XOR.

For a nonzero input difference $\Delta x$, an S-box is called \emph{active}.
The difference distribution table (DDT) is
\begin{equation}\label{eq:ddt}
 D[\alpha,\beta]
 =\bigl|\{u\in\{0,\ldots,15\}: S(u)\oplus S(u\oplus\alpha)=\beta\}\bigr|,
 \qquad p(\alpha,\beta)=\frac{D[\alpha,\beta]}{16}.
\end{equation}
The computed DDT contains 159 zero entries, 72 entries equal to 2,
24 entries equal to 4, and one entry equal to 16. Thus 97 of the 256
input/output difference pairs are allowed. For an active S-box,
the largest single-transition probability is $4/16=2^{-2}$.

\subsection{Characteristic probability convention}
A differential \emph{characteristic} specifies the input and output difference
of every S-box along the chosen four-round path. Its standard DDT-product weight is
\begin{equation}\label{eq:trail-prob}
 P_{\mathrm{trail}}=\prod_{r=0}^{3}\prod_{c=0}^{15}
 \frac{D[\Delta x_{r,c},\Delta y_{r,c}]}{16}.
\end{equation}
This is a product of local differential transition probabilities; it is not,
without additional assumptions, a measured probability of the entire specified
path for an individual fixed key. It is also distinct from a \emph{differential}
probability obtained by summing over multiple characteristics with matching
external differences.

\section{Exact MILP Formulation}
\subsection{S-box activity indicators}
Let $x_0,x_1,x_2,x_3$ be the four input difference bits of a bijective S-box;
let $y_0,y_1,y_2,y_3$ denote its output difference bits; and introduce an
activity indicator $A\in\{0,1\}$. The seven activity and nonzero-propagation
inequalities used in \cite[Sec.~4.1]{ilter2023} are
\begin{align}
 x_i-A&\leq 0 &&(i=0,1,2,3),\label{eq:act}\\
 \sum_{i=0}^{3}x_i-A&\geq 0,\label{eq:act-lower}\\
 4\sum_{i=0}^{3}x_i-\sum_{i=0}^{3}y_i&\geq0,\\
 4\sum_{i=0}^{3}y_i-\sum_{i=0}^{3}x_i&\geq0.
\end{align}
These ensure that $A=1$ precisely when the input difference is nonzero and
that a nonzero input difference cannot produce a zero output difference (or
vice versa). There are $7\cdot64=448$ such inequalities in the four-round model.
They encode activity but do not, by themselves, capture the complete DDT.

\subsection{Exact S-box transitions and probability costs}
To encode exact S-box transitions and probabilities, we follow the
H-representation and set-cover approach of Sasaki and Todo
\cite{sasaki2017}, as used in \cite[Sec.~4.1]{ilter2023}.
Every S-box transition is represented by a ten-bit vector
\[
 E=(x_3,x_2,x_1,x_0,y_3,y_2,y_1,y_0,P_0,P_1).
\]
The two probability bits encode $-\log_2 p$ as follows:
\begin{table}[H]
 \centering
 \caption{Exact probability encoding for allowed Midori-64 S-box transitions.}
 \begin{tabular}{ccccc}
 \toprule
 DDT count & Local probability & $P_0$ & $P_1$ & $P_0+2P_1$\\
 \midrule
 $16$ & $1$ & 0 & 0 & 0\\
 $4$ & $2^{-2}$ & 0 & 1 & 2\\
 $2$ & $2^{-3}$ & 1 & 1 & 3\\
 \bottomrule
 \end{tabular}
\end{table}
The underlying SageMath construction forms the convex hull of all 97 allowed
binary vectors, computes its H-representation, and uses a set-cover MILP to
select inequalities excluding every disallowed binary vector. The resulting
system consists of \textbf{19 inequalities per S-box}. The precise 19 rows
are provided in \texttt{selected\_inequalities.txt} and are loaded unchanged
by the model generator. An exhaustive verification over all $2^{10}=1024$
binary vectors confirms that \emph{exactly} the 97 DDT-consistent vectors
are accepted; all 927 other assignments are rejected. Over 64 S-box instances,
this contributes $19\cdot64=1,216$ inequalities.

The seven activity inequalities and the 19 exact-DDT inequalities
are included together in \emph{both} full-round optimization models. This
shared S-box representation is central to methodological consistency.

\subsection{ShuffleCell permutation constraints}
For the state serialization used here, the ShuffleCell output at position $c$
reads the input at position $\pi(c)$, where
\begin{equation}\label{eq:shuffle}
\pi=(0,10,5,15,14,4,11,1,9,3,12,6,7,13,2,8).
\end{equation}
Let $y_{r,c,b}$ be bit $b$ of the SubCell output difference at round $r$ and
$z_{r,c,b}$ the corresponding ShuffleCell output bit. The permutation is
modeled exactly by the equality
\begin{equation}\label{eq:sh-eq}
 z_{r,c,b}=y_{r,\pi(c),b},\qquad
 r=0,1,2;\quad c=0,\ldots,15;\quad b=0,\ldots,3.
\end{equation}
Three layers of 64 bit equalities produce $192$ constraints, in the manner of
\cite[Sec.~4.2]{ilter2023}.

\subsection{MixColumn and XOR parity constraints}
Each column is represented by four consecutive cells. With column input
$(z_0,z_1,z_2,z_3)$, the Midori MixColumn transformation is
\begin{equation}\label{eq:mix}
\begin{pmatrix}w_0\\w_1\\w_2\\w_3\end{pmatrix}=
\begin{pmatrix}
0&1&1&1\\1&0&1&1\\1&1&0&1\\1&1&1&0
\end{pmatrix}
\begin{pmatrix}z_0\\z_1\\z_2\\z_3\end{pmatrix}
\quad\text{over }\mathbb F_2^4.
\end{equation}
In particular, $w_0=z_1\oplus z_2\oplus z_3$. For three input bits
$a,b,c$ and output bit $d=a\oplus b\oplus c$, the parity is encoded by
\begin{equation}\label{eq:xor}
 a+b+c+d=4u-2v,\qquad u,v\in\{0,1\}.
\end{equation}
The output of each inter-round MixColumn becomes the next SubCell input
$x_{r+1,c,b}$. Using \eqref{eq:xor} for every bit gives $192$ parity
equalities and $384$ auxiliary binary variables. Exhaustive testing of
all $16$ assignments $(a,b,c,d)$ confirms that the equation is feasible
in the dummy variables if and only if $d=a\oplus b\oplus c$.

\section{Four-Round Model and Experimental Setup}
\subsection{A common feasible region for both objectives}
Let $A_{r,c}$ denote S-box activity and $P_{0,r,c},P_{1,r,c}$ its probability bits.
We use the same four-round constraints for both experiments and require a
nontrivial initial difference:
\begin{equation}\label{eq:nonzero}
 \sum_{c=0}^{15}\sum_{b=0}^{3} x_{0,c,b}\geq1.
\end{equation}
Only the objective changes:
\begin{align}
 \text{Active S-box count: } &\min N_{\rm active}
 =\min\sum_{r=0}^{3}\sum_{c=0}^{15}A_{r,c},\label{eq:obj-act}\\
 \text{Characteristic cost: } &\min C
 =\min\sum_{r=0}^{3}\sum_{c=0}^{15}
 \bigl(P_{0,r,c}+2P_{1,r,c}\bigr)
 =\min\bigl(-\log_2P_{\rm trail}\bigr).\label{eq:obj-prob}
\end{align}

\begin{table}[H]
 \centering
 \caption{Size of each four-round MILP before presolve.}
 \begin{tabular}{lr}
 \toprule
 Constraint category & Number\\
 \midrule
 S-box activity inequalities ($64\times7$) & 448\\
 S-box H-representation inequalities ($64\times19$) & 1,216\\
 ShuffleCell bit equalities ($3\times64$) & 192\\
 MixColumn XOR equalities ($3\times64$) & 192\\
 Nonzero initial difference & 1\\
 \midrule
 \textbf{Total constraints} & \textbf{2,049}\\
 \textbf{Total binary variables} & \textbf{1,280}\\
 \bottomrule
 \end{tabular}
\end{table}
The $1,280$ variables consist of 512 SubCell input/output difference bits,
64 activity indicators, 128 probability bits, 192 ShuffleCell bits, and
384 XOR dummy bits. The LP generator writes separate files for the two
objectives while keeping their \emph{constraint systems identical}.

\subsection{Validation and solver environment}
The component tests check the entire ten-bit H-representation truth table,
the XOR parity encoding, the agreement of the two models' constraint systems,
and randomly constructed valid differential trails. Fourteen automated tests
pass, including a replay of the published experiment records. Independently, the saved Gurobi activity solution was checked against
all 2,049 inequalities with zero violations; the probability trail was checked
against a separately implemented Midori DDT and linear layers.

Optimization used Gurobi 13.0.3 on Windows 11 with an AMD Ryzen AI 7 350 CPU.
The final activity-count run used eight threads, \texttt{MIPFocus=3},
\texttt{MIPGap=0}, and a six-hour time limit. A previously obtained feasible
probability-trail solution with objective 16 was supplied as a MIP start; the
new run \emph{certified optimality}, rather than discovering this incumbent. The solver terminated
well before the time limit.

\section{Certified Minimum Number of Active S-boxes}
Gurobi \emph{certified} a minimum of 16 active S-boxes for the
four-round MILP described in Section~4. The run loaded a feasible 16-active
solution as a MIP start (recorded as \texttt{Loaded user MIP start with
objective 16}). The solver's contribution was to close the lower bound to
16 and certify global optimality within the model; it did not discover the
initial incumbent during this run. The final solver information is
summarized in Table~\ref{tab:active-res}.
\begin{table}[H]
 \centering
 \caption{Certified four-round active-S-box result (Gurobi 13.0.3).}\label{tab:active-res}
 \begin{tabular}{lr}
 \toprule
 Solver status & \texttt{OPTIMAL}\\
 Best feasible objective (incumbent) & 16\\
 Best objective bound & 16\\
 Relative optimality gap & $0\%$\\
 Branch-and-bound nodes explored & 60,345\\
 Runtime & $794.455$ s ($\approx13$ min $14$ s)\\
 Model size & 2,049 constraints; 1,280 binary variables\\
 \bottomrule
 \end{tabular}
\end{table}
Because the solver's certified lower bound matches the feasible incumbent,
\begin{equation}\label{eq:active-result}
 \boxed{\min N_{\rm active}=16}
\end{equation}
within this exact four-round model (subject to standard numerical tolerances).
This independently reproduces the four-round value reported for Midori64 by
its designers \cite{banik2015}. It is \emph{not} a newly discovered
active-S-box lower bound for the cipher.

An optimal solution has the following active cell indices, numbered from 0
through 15:
\begin{table}[H]
 \centering
 \caption{Activity pattern of the saved optimal solution.}
 \begin{tabular}{clc}
 \toprule
 Round & Active cell indices & Count\\
 \midrule
 1 & $0,10,15$ & 3\\
 2 & $2$ & 1\\
 3 & $12,13,15$ & 3\\
 4 & $0,1,2,8,9,11,12,14,15$ & 9\\
 \midrule
 Total & & \textbf{16}\\
 \bottomrule
 \end{tabular}
\end{table}

\subsection{The wide-trail bound and independent word-level validation}
The theoretical explanation predates this implementation. Midori64's
MixColumn has differential branch number $B=4$, and its ShuffleCell is an
optimal cell permutation for diffusion. Under these conditions, the
four-round propagation theorem of Daemen and Rijmen
\cite[Ch.~9, Sec.~9.5]{daemen2002} gives at least $B^2=16$ active
S-boxes across four rounds. The original Midori paper reports this value
in its active-S-box comparison (Table~8) \cite{banik2015}.
The following dynamic program provides a separate, directly executable
verification of the bound for Midori's \emph{actual} column transformation
and cell permutation, without invoking a floating-point MILP solver.

A separate verification does not require MILP or numerical solver tolerances.
Represent an active-cell pattern by a nonzero binary word
$u\in\{0,1\}^{16}$ and let $w(u)$ be its Hamming weight. For each
four-cell input pattern to MixColumn, enumerate the possible nonzero output
patterns by applying the exact linear transformation to all $16^4$
nibble vectors. Let $T$ denote the induced, four-column relation, composed
with the ShuffleCell permutation. Then the dynamic program
\[
 F_1(u)=\begin{cases}w(u),&u\ne0,\\+\infty,&u=0,\end{cases}
 \qquad
 F_{r+1}(v)=w(v)+\min_{u:\,u\,T\,v}F_r(u)
\]
minimizes the sum of active-cell counts over the \emph{relaxed} activity model.
Because every concrete bit-level characteristic induces a permissible
sequence of activity patterns, its active-S-box count is at least
$\min_v F_4(v)$. The independently implemented program
\texttt{independent\_checks.py} produces, for one through eight rounds,
\[
 (1,4,7,16,23,30,35,38).
\]
In particular, $\min_v F_4(v)=16$, and our explicit bit-level
16-active characteristic attains this lower bound. This complements the Gurobi calculation with a combinatorial,
solver-independent lower bound. Together with the explicit 16-active
bit-level characteristic, it establishes equality in the stated four-round
setting. This reproduces a known diffusion result; it is not a new
cryptanalytic bound.

\subsection{An equivalent compact LP}
The 192 ShuffleCell bit equalities can be eliminated by substituting
$z_{r,c,b}=y_{r,\Pi(c),b}$ directly into the MixColumn equations.
The resulting LP has 1,857 constraints and 1,088 binary variables. An
automated equivalence audit checks that the archived full-model optimum
projects onto a feasible compact-model assignment of objective 16.
The compact construction is designed to fit Gurobi's standard 2,000-row,
2,000-variable size-limited license; the archived \texttt{OPTIMAL}
solver log is nevertheless for the original full model.

\section{Best Four-Round Differential Characteristic}
A separate probability-minimizing Gurobi run found the characteristic in
Table~\ref{tab:trail}. The direct probability optimization was interrupted
after $2,289.24$ seconds: its best cost was 32, its branch-and-bound lower
bound was 18, and its reported gap was $43.75\%$. Therefore, \emph{that
run alone did not certify probability optimality}. The certified
active-S-box result supplies the missing bound on the DDT-product weight.

\begin{table}[H]
 \centering
 \caption{Four-round Midori-64 differential characteristic; states are 16 hex nibbles in column-major order.}\label{tab:trail}
 \footnotesize
 \begin{tabular}{c@{\hspace{5pt}}l@{\hspace{7pt}}l@{\hspace{6pt}}cc}
 \toprule
 Round & SubCell input $\Delta x$ & SubCell output $\Delta y$ & Active & $p_r$\\
 \midrule
 1 & \texttt{2000000000200002} & \texttt{4000000000400004} & 3 & $2^{-6}$\\
 2 & \texttt{0040000000000000} & \texttt{0020000000000000} & 1 & $2^{-2}$\\
 3 & \texttt{0000000000002202} & \texttt{0000000000004404} & 3 & $2^{-6}$\\
 4 & \texttt{4440000044044044} & \texttt{2220000022022022} & 9 & $2^{-18}$\\
 \midrule
 Total & & & \textbf{16} & $\boldsymbol{2^{-32}}$\\
 \bottomrule
 \end{tabular}
\end{table}
Each of the 16 active transitions in this path has a DDT count of 4 and
therefore probability $2^{-2}$. Inactive S-boxes contribute probability 1.
Consequently,
\begin{equation}\label{eq:prob-witness}
 P_{\rm trail}=2^{-6}\cdot 2^{-2}\cdot 2^{-6}\cdot 2^{-18}
 =\boxed{2^{-32}}, \qquad C=32.
\end{equation}
The reported SubCell outputs propagate exactly to the following round's
inputs under ShuffleCell and MixColumn. This was independently checked using
a reference implementation rather than trusting the MILP encoding alone.

\begin{proposition}[Characteristic optimality in the specified MILP model]\label{prop:opt}
Within the four-round model of Sections 2--4, the maximal DDT-product
probability of a nontrivial differential characteristic is $2^{-32}$.
\end{proposition}
\begin{proof}
The active-S-box minimization certifies that every feasible nontrivial four-round characteristic has
at least 16 active S-boxes. The maximal DDT probability for each active
Midori-64 S-box transition is $2^{-2}$, while an inactive transition has
probability 1. Thus, for every such characteristic,
\[
 P_{\rm trail}\leq(2^{-2})^{N_{\rm active}}\leq2^{-32}.
\]
The characteristic in Table~\ref{tab:trail} achieves $2^{-32}$ by
\eqref{eq:prob-witness}, so the bound is attained.
\end{proof}
This proposition establishes optimality for \emph{single characteristics in
the modeled four-round construction}. It does not assert an optimum for the
full 16-round cipher or for differential probabilities aggregated over trails.

\section{Independent Software Verification}
\subsection{Layer-by-layer validation}
A separately implemented four-round Midori core recomputed the DDT and
checked all 64 local nibble transitions. The specified differences are
valid and the four active-S-box counts are $(3,1,3,9)$. Independent
ShuffleCell/MixColumn operations map each of the first three SubCell output
differences exactly to the next SubCell input difference. These checks pass.
However, locally valid S-box transitions do not on their own demonstrate that
one plaintext pair realizes the entire multi-round path.

\subsection{Constructive plaintext-pair verification}
A second MILP held the four-round difference trail fixed and searched for a
128-bit master key and a pair of plaintexts realizing every internal transition.
The reduced Midori-64 implementation uses $K=K_0\Vert K_1$, the whitening key
$WK=K_0\oplus K_1$, three round-key additions with
\begin{equation}\label{eq:key}
 RK_i=K_{i\bmod 2}\oplus\alpha_i\qquad(i=0,1,2),
\end{equation}
and initial and final whitening. The fixed $\alpha_i$ are the first three
Midori-64 round constants in the model's state serialization. The configuration
has three full rounds and a final SubCell, hence \emph{four SubCell rounds in
total}; it is not a full 16-round encryption.

The search returned the master key
\[
 K=\texttt{00000101000100016177166671066660},
\]
whitening key $WK=\texttt{6177176771076661}$, and round keys
\begin{align*}
 RK_0&=\texttt{0001000010100010},\\
 RK_1&=\texttt{6066066660066660},\\
 RK_2&=\texttt{1010000100100100}.
\end{align*}
With this master key, the following explicit plaintexts and ciphertexts were
obtained:
\begin{table}[H]
 \centering
 \caption{Constructive four-round witness under the displayed master key.}\label{tab:pair}
 \footnotesize
 \begin{tabular}{lll}
 \toprule
 & First text & Second text\\
 \midrule
 Plaintext & \texttt{ADBBDBABBDCBAAAD} & \texttt{8DBBDBABBDEBAAAF}\\
 Ciphertext & \texttt{ADBDDBADBDABACAD} & \texttt{8F9DDBAD9FA98C8F}\\
 \bottomrule
 \end{tabular}
\end{table}
In particular,
\begin{align*}
 \Delta P&=\texttt{2000000000200002},\\
 \Delta C&=\texttt{2220000022022022}.
\end{align*}
Independent execution of both plaintexts through the reduced-round
implementation reproduces \emph{every} SubCell input/output difference in
Table~\ref{tab:trail}, not merely the external differences. The result is an
explicit end-to-end witness of feasibility for one standard-schedule master
key in the specified reduced four-round variant. For comparison, this exact
trail was infeasible in a separate witness search where the master key was
fixed to all zeros; realizability of an exact path can depend on the key.

The witness MILP had an identically zero objective. Its solver status
\texttt{OPTIMAL} means that a feasible pair was found; this must not be
confused with optimization of the differential characteristic. Finally, one
witness does not measure a $2^{-32}$ empirical occurrence rate for the
fixed-key encryption permutation; a separate conditional sampling
experiment for this key is reported in Section~\ref{sec:fixed-key}.

\section{Exploratory Fixed-Key Characteristic Behaviour}
\label{sec:fixed-key}
The $2^{-32}$ weight computed from local DDT entries is not the
fixed-key probability of the specified path. In particular, both active
transitions used in the witness involve affine conforming-input sets:
\[
 2\to4:\ \{\mathtt C,\mathtt D,\mathtt E,\mathtt F\},\qquad
 4\to2:\ \{0,2,4,6\}.
\]
The $2\to4$ conforming outputs lie in the XOR subgroup
$E=\{0,2,4,6\}$. At the sole active S-box of round two,
\[
 x_{1,2}=S(x_{0,0})\oplus S(x_{0,10})\oplus
 S(x_{0,15})\oplus RK_{0,2}.
\]
Consequently, the round-two transition is compatible with all conforming
first-round triples if $RK_{0,2}\in E$, and incompatible with all such
triples otherwise. A uniformly chosen master key satisfies this necessary
condition with probability $1/4$. This is not a complete characterization
of the keys for which all four rounds can follow the path.

A separate C implementation performs conditional sampling for the
specific verified master key: it chooses first-round active inputs
uniformly among the conforming triples and samples the other nibbles.
The exact conditioning probability $2^{-6}$ is incorporated when
estimating the full-path probability. Two runs with $10^8$ conditional
samples each yielded 1,541 and 1,543 hits, corresponding to
estimated probabilities approximately $2^{-21.99}$ and $2^{-21.98}$.
Thus the observed fixed-key path probability is near $2^{-22}$,
well above the DDT-product $2^{-32}$. These are statistical estimates,
not an analytic probability or a claim about the fraction of compatible
master keys. We interpret them in the established fixed-key and
plateau-characteristic literature, in particular Canteaut and Fruchon
\cite{canteaut2025}. The underlying phenomenon is not claimed as new.

\section{Discussion, Limitations, and Conclusion}
This reproducible implementation recovers the known four-round Midori-64
minimum of 16 active S-boxes by two independent methods: Gurobi's
\texttt{OPTIMAL} certificate for the exact bit-level MILP, and a
solver-free word-level dynamic program. The latter is consistent with the
classical wide-trail four-round propagation bound. Because the maximum
probability of a nonzero Midori S-box transition is $2^{-2}$, every
nontrivial four-round characteristic has a DDT-product of at most
$2^{-32}$. The explicit characteristic in Table~\ref{tab:trail} attains
this bound. Its realizability is confirmed by a master key and plaintext
pair under the standard key schedule in a reduced four-round Midori-64
implementation.

The DDT-product is not the actual occurrence probability of the same
path under every fixed master key. Conditional sampling for the displayed
master key estimates a path probability near $2^{-22}$, whereas the
round-two algebraic condition in Section~\ref{sec:fixed-key} already
excludes three quarters of uniformly chosen keys. This illustrates
fixed-key/plateau behaviour studied in the existing literature, rather
than a new attack or security result. Deriving the complete set of
compatible-key conditions and an exact fixed-key probability is a
natural direction for further investigation. The present contribution
is the transparent modeling, exhaustive component tests, independent
lower-bound calculation, saved optimization record, and constructive
software verification of a known four-round benchmark.

\appendix
\section{Reproducibility and Saved Artifacts}
The model builder (\texttt{build\_models.py}) and solver
(\texttt{solve.py}) share the saved inequalities. Software verification
uses \texttt{midori\_reference.py} and \texttt{verify\_trail.py}.
Additional scripts perform focused active-S-box optimization, an independent
constraint audit, and a standard-key witness search. In the repository root, the principal commands are:
\begin{lstlisting}
python build_models.py
python -u -m unittest discover -s tests -v
python -u overnight_midori_active.py --hours 6
python -u solve.py --objective probability --time-limit 7200
python independent_checks.py
python compact_models.py
python -u verify_trail.py --report results/probability_trail.json --samples 0
python -u find_midori_standard_key_pair.py --report results/probability_trail.json --search-master-key --time-limit 120
\end{lstlisting}
The retained experiment records and reproducibility tools provide the audit trail:
\begin{itemize}
\item \texttt{results/active\_optimality.json} --- certified 16/16 objective and bound, status \texttt{OPTIMAL};
\item \texttt{results/active\_solver.log} --- full branch-and-bound record;
\item \texttt{results/active\_solution.sol} --- the complete 1,280-variable assignment;
\item \texttt{results/probability\_trail.json} --- objective 32 and all four rounds' difference states;
\item \texttt{results/standard\_key\_witness.json} --- concrete master key, plaintexts, ciphertexts, and witness check;
\item \texttt{results/fixed\_key\_sampling.json} --- conditional Monte Carlo settings, seeds, and observed hits.
\end{itemize}
The probability-only run was interrupted and does not have solver-certified
optimality. The characteristic optimum follows instead from the certified active-S-box bound and
Proposition~\ref{prop:opt}. Numerical solver certification is subject to
Gurobi tolerances, as in standard MILP practice. A word-level dynamic program independently
provides the same 16-active lower bound. The archived standard-key witness is separately replayed against a
test-vector-validated Midori implementation. The repository also provides a
small feasibility MILP to obtain a fresh witness, which may differ from
the recorded one.

\begin{thebibliography}{9}
\bibitem{banik2015}
S. Banik, A. Bogdanov, T. Isobe, K. Shibutani, H. Hiwatari, T. Akishita,
and F. Regazzoni.
\newblock Midori: A block cipher for low energy.
\newblock In \emph{Advances in Cryptology---ASIACRYPT 2015}, Part II,
LNCS 9453, pp. 411--436, Springer, 2015.
\newblock \url{https://eprint.iacr.org/2015/1142}.

\bibitem{ilter2023}
M. B. \c{I}lter and A. A. Sel\c{c}uk.
\newblock MILP-aided cryptanalysis of the FUTURE block cipher.
\newblock In \emph{Innovative Security Solutions for Information Technology
and Communications (SecITC 2022)}, LNCS 13809, pp. 153--167,
Springer, 2023.
\newblock \url{https://doi.org/10.1007/978-3-031-32636-3_9}.
\newblock Preprint: \url{https://eprint.iacr.org/2022/1398}.

\bibitem{sasaki2017}
Y. Sasaki and Y. Todo.
\newblock New algorithm for modeling S-box in MILP based differential and
division trail search.
\newblock In \emph{Innovative Security Solutions for Information Technology
and Communications (SecITC 2017)}, LNCS 10543, pp. 150--165, Springer, 2017.
\newblock \url{https://doi.org/10.1007/978-3-319-69284-5_11}.
\bibitem{mouha2012}
N. Mouha, Q. Wang, D. Gu, and B. Preneel.
\newblock Differential and linear cryptanalysis using mixed-integer linear programming.
\newblock In \emph{Inscrypt 2011}, LNCS 7537, pp. 57--76, Springer, 2012.
\newblock \url{https://doi.org/10.1007/978-3-642-34704-7_5}.

\bibitem{sun2014}
S. Sun, L. Hu, P. Wang, K. Qiao, X. Ma, and L. Song.
\newblock Automatic security evaluation and (related-key) differential characteristic search.
\newblock In \emph{ASIACRYPT 2014}, LNCS 8873, pp. 158--178, Springer, 2014.
\newblock \url{https://eprint.iacr.org/2013/676}.

\bibitem{daemen2002}
J. Daemen and V. Rijmen.
\newblock \emph{The Design of Rijndael: AES---The Advanced Encryption Standard}.
\newblock Springer, 2002, Chapter 9, Section 9.5 (Four-Round Propagation Theorem).
\newblock \url{https://doi.org/10.1007/978-3-662-04722-4}.

\bibitem{canteaut2025}
A. Canteaut and M. Fruchon.
\newblock Understanding unexpected fixed-key differential behaviours: How to avoid major weaknesses in lightweight designs.
\newblock \emph{ASIACRYPT 2025}. Extended version: \url{https://eprint.iacr.org/2025/1553}.
\end{thebibliography}

\end{document}
